
par (i) \(\Rightarrow\) (iv) prenons une surjection \(L={O}_{X_T}(-m)^{\oplus M}\longrightarrow E\). Quitte à couper \(T\) en morceaux, de sorte que \(E\) devienne plat sur \(T\), la formation de la suite exacte \(0\longrightarrow I\xrightarrow{u}L\longrightarrow E\longrightarrow0\) commute au passage aux fibres. Par (i) \(\Rightarrow\) (iii) et (i) \(\Rightarrow\) (iv), prenons \(m_1\) assez grand pour que tous les \(\operatorname{Ext}^1({O}_{X_K}(-m_1),I_K)=H^1(I_K(m_1))\) soient nuls et qu'il existe des surjections \(L_1={O}_{X_T}(-m_1)^{\oplus M_1}\longrightarrow E'\) et \(\alpha:L_2={O}_{X_T}(-m_1)^{\oplus M_2}\longrightarrow I\). Alors les applications \(\operatorname{Hom}(L_{1,K},L_K)\longrightarrow \operatorname{Hom}(L_{1,K},E_K)\) sont surjectives. Soit \(\beta:E'_K\to E_K\) un homomorphisme et soit \(\gamma\) la composition \(L'_K\to E'_K\to E_K\); alors \(\gamma\) provient d'un homomorphisme \(\delta:L_{1,K}\to L_K\), et $(\delta,u_K\circ\alpha):L_{1,K}\oplus L_{2,K}\longrightarrow L_K$ et \(\beta:E'_K\to E_K\) ont le même conoyau. Par conséquent, \(\mathcal F\) est contenue dans la famille des classes des conoyaux d'homomorphismes de la forme \({O}_{X_K}(-m_1)^{\oplus(M_1+M_2)}\longrightarrow {O}_{X_K}(-m)^{\oplus M}\).

Toujours quitte à couper \(T\) en morceaux, on peut supposer en plus que \(X_T\) (donc \(\underline{\Hom}_{{O}_{X_T}}(E',E)\)) et tous les \(R^q(f_T)_*\underline{\Hom}_{{O}_{X_T}}(E',E)\) (où \(f_T:X_T\to T\) est le morphisme structural) sont plats sur \(T\), afin que la formation de \(G=(f_T)_*\underline{\Hom}_{{O}_{X_T}}(E',E)\) commute au changement de base quelconque \(T'\to T\) (EGA III 6.9.9.2). Soit $R=\mathbb V(G^*)$, alors \(X_R\) porte canoniquement une suite exacte ``universelle''
\[
E'_R\longrightarrow E_R\longrightarrow F\longrightarrow0
\]
et \(F\) limite \(\mathcal F\).

\subsection*{2. Plusieurs lemmes techniques}
\textbf{2.1.} Soient \(k\) un corps et \(X\) un \(k\)-schéma propre. Soient \(L_1,\ldots,L_s\) des faisceaux inversibles sur \(X\), \(c_1(L_i)\) la classe de Chern de \(L_i\) et \(Y\) un sous-schéma fermé de \(X\) de dimension \(\le s\). Rappelons qu'on désigne par
